% ##############################################################################
\section{Contracts and notation}
\label{sec:contracts_and_notation}

This section fixes the notation used throughout the rest of the paper: the
capability tier set, the task unit, the contract, and the bid/ask tuples
submitted to \NoesisMarket.

\begin{definition}[Capability tiers]
\label{def:capability_tiers}
Let
$$
  \calC = \{ c_1, c_2, \ldots, c_K \}
$$
be a finite, totally ordered set of \emph{capability tiers}, with
$c_1 \prec c_2 \prec \cdots \prec c_K$, where $\prec$ denotes ``strictly
less capable than''. A tier is a discrete, provider-agnostic label (e.g.,
\emph{cheap}, \emph{medium}, \emph{frontier-class}) that a buyer can
require and a seller can offer, rather than a continuous benchmark score.
\end{definition}

Tiers are deliberately coarse. A continuous capability score would be more
informative but is harder for a seller to credibly commit to ex ante and
harder for a buyer to reason about ex post; Section~\ref{sec:open_questions}
returns to this trade-off.

\begin{definition}[Task unit]
\label{def:task}
A \emph{task} is an abstract, normalized unit of inference work, denoted
$u \in \calU$, used as the unit of quantity in a contract. Throughout the
formalization, quantities are expressed in tasks without committing to a
specific definition of $\calU$.
\end{definition}

Definition~\ref{def:task} is deliberately left abstract. Candidate
definitions include a raw token count, a wall-clock compute duration, or a
benchmark-normalized task-equivalent that adjusts for the fact that a
token from a frontier model and a token from a small model are not
comparable units of work. This choice affects cross-provider
comparability directly and is treated as an open research question
(Section~\ref{sec:open_questions}) rather than resolved here; the
remainder of the formalization only requires that $\calU$ support addition
and a total order, both of which hold for any of the candidate
definitions above.

\begin{definition}[Intelligence contract]
\label{def:contract}
An \emph{intelligence contract} is a tuple
$$
  \kappa = ( \ntasksattr, \clevelattr, \lmaxattr, \rminattr, \priceattr )
$$
where
\begin{itemize}
  \item $\ntasksattr \in \bbN$ is the number of tasks covered by the
    contract,
  \item $\clevelattr \in \calC$ is the required capability tier,
  \item $\lmaxattr \in \bbR_+$ is the maximum latency per task,
  \item $\rminattr \in [0, 1]$ is the minimum reliability, i.e., the
    minimum fraction of tasks that must meet both $\clevelattr$ and
    $\lmaxattr$,
  \item $\priceattr \in \bbR_+$ is the price of the contract.
\end{itemize}
\end{definition}

The quantities are ordered to read naturally: 

\emph{``deliver $\ntasksattr$ tasks at capability $\clevelattr$ or better,
within $\lmaxattr$ latency, with at least $\rminattr$ reliability, for a total
price of $\priceattr$.''}

\begin{example}[Intelligence contract]
\label{ex:contract}
The contract
$$
  \kappa_1 = (100,\ \texttt{frontier},\ 2\text{s},\ 0.999,\ P)
$$
reads: ``deliver 100 tasks at frontier-class capability or better,
within 2 seconds of latency, with at least 99.9\% reliability, for a total
price of $P$.''
\end{example}

% ==============================================================================
\subsection{Bids and asks}
\label{sec:bids_asks}

A contract $\kappa$ is the outcome of \NoesisMarket{} matching a bid
against one or more compatible asks. This paper instantiates matching as
a periodic call auction (Section~\ref{sec:noesis_market}), but the
contract definition itself does not depend on how bids and asks are
matched: a continuous double auction or a posted-price mechanism
(Section~\ref{sec:modularity}) would produce a contract of the same
shape. Bids and asks share the contract's attributes, plus the
buy/sell-specific role of the price attribute as a limit rather than a
fixed value.

\begin{definition}[Bid]
\label{def:bid}
A \emph{bid} is a tuple
$$
  \beta = (n_\beta, c_\beta, \ell_\beta, r_\beta, p_\beta)
$$
submitted by a buyer, where
\begin{itemize}
  \item $n_\beta$ is the requested number of tasks,
  \item $c_\beta \in \calC$ is the minimum required capability tier,
  \item $\ell_\beta$ is the maximum acceptable latency,
  \item $r_\beta$ is the minimum acceptable reliability,
  \item $p_\beta$ is the maximum price per task the buyer is willing to
    pay.
\end{itemize}
\end{definition}

\begin{definition}[Ask]
\label{def:ask}
An \emph{ask} is a tuple
$$
  \alpha = (n_\alpha, c_\alpha, \ell_\alpha, r_\alpha, p_\alpha)
$$
submitted by a seller, where
\begin{itemize}
  \item $n_\alpha$ is the offered number of tasks,
  \item $c_\alpha \in \calC$ is the offered capability tier,
  \item $\ell_\alpha$ is the seller's typical (self-reported or previously
    measured) latency,
  \item $r_\alpha$ is the seller's typical reliability,
  \item $p_\alpha$ is the minimum price per task the seller is willing to
    accept.
\end{itemize}
\end{definition}

Unlike a bid, an ask's capability, latency, and reliability attributes are
descriptive of what the seller can deliver rather than a constraint on
what the seller will accept; only $p_\alpha$ is a limit in the auction
sense. Section~\ref{sec:compatibility} defines when a bid and an ask are
compatible, i.e., when the ask's descriptive attributes satisfy the bid's
required attributes.

Figure~\ref{fig:contract_lifecycle} summarizes the resulting lifecycle: a
bid (Definition~\ref{def:bid}) and one or more compatible asks
(Definition~\ref{def:ask}) are matched by \NoesisMarket{} into a single
intelligence contract (Definition~\ref{def:contract}), which is then
dispatched to \NoesisServer{} for execution.

% rendered_images:begin
% ```graphviz
% digraph ContractLifecycle {
%   bgcolor="transparent";
%   pad="0.12";
%   splines=spline;
%   nodesep=0.40;
%   ranksep=0.55;
%   rankdir=LR;
% 
%   node [shape=box,
%         style="rounded,filled",
%         fillcolor="#FFFFFF",
%         color="#9AA9B8",
%         penwidth=1.8,
%         fontname="Helvetica",
%         fontcolor="#243B53",
%         fontsize=11,
%         margin="0.20,0.14",
%         height=0.50];
% 
%   edge [color="#A3B1C0",
%         penwidth=1.3,
%         arrowhead=vee,
%         arrowsize=0.75,
%         fontname="Helvetica",
%         fontsize=9,
%         fontcolor="#7B8794"];
% 
%   Bid [label=<<b>Bid</b><br/><font point-size="9" color="#A34F6C">Buyer</font>>,
%        fillcolor="#F6E1E8", color="#D98CA8", fontcolor="#6B2A44"];
%   Ask [label=<<b>Ask(s)</b><br/><font point-size="9" color="#946A2E">Seller(s)</font>>,
%        fillcolor="#FBEBD4", color="#D9A85F", fontcolor="#6B4517"];
%   Matching [label=<<b>NoesisMarket</b><br/><font point-size="9" color="#41719C">Matches bid against asks</font>>,
%             fillcolor="#D3E3F3", color="#7CA6CE", fontcolor="#1F4E79"];
%   Contract [label=<<b>Intelligence contract</b><br/><font point-size="9" color="#2E7C74">(n, c, L_max, R_min, P)</font>>,
%             fillcolor="#D4EEEB", color="#7FC2B9", fontcolor="#1F5A54"];
%   Server [label=<<b>NoesisServer</b><br/><font point-size="9" color="#41719C">Executes contract</font>>,
%           fillcolor="#D3E3F3", color="#7CA6CE", fontcolor="#1F4E79"];
% 
%   Bid -> Matching [label="bid"];
%   Ask -> Matching [label="compatible ask(s)"];
%   Matching -> Contract [label="clears"];
%   Contract -> Server [label="dispatched"];
% }
% ```
% label=fig:contract_lifecycle
% caption=The contract lifecycle: a bid and one or more compatible asks are matched by \NoesisMarket{} into an intelligence contract, dispatched to \NoesisServer{} for execution.
% rendered_images:end
% render_images:begin
\begin{figure}[H]
  \includegraphics[width=\linewidth]{03_contracts_and_notation.tex.figs/03_contracts_and_notation.1.png}
  \caption{The contract lifecycle: a bid and one or more compatible asks are matched by \NoesisMarket{} into an intelligence contract, dispatched to \NoesisServer{} for execution.}
  \label{fig:contract_lifecycle}
\end{figure}
% render_images:end
